\documentclass{article}
\usepackage[T1]{fontenc}
\usepackage{lmodern}
\usepackage{graphicx}
\usepackage{amsmath,amssymb}
\usepackage[a4paper,margin=2cm]{geometry}
\usepackage[absolute,overlay]{textpos}
\setlength{\TPHorizModule}{1mm}
\setlength{\TPVertModule}{1mm}
\usepackage[indonesian]{babel}
\usepackage{hyperref}
\setlength{\emergencystretch}{2em}
% unit: Halaman 43

\begin{document}

% === Halaman 43 (llm) ===

\textbf{Contoh *):} Misalkan kurva eliptik yang dipilih adalah $y^2 = x^3 + 2x + 1 \pmod 5$. Himpunan titik-titik pada kurva eliptik adalah $\{\{(0, 1), (1, 3), (3, 3), (3, 2), (1, 2), (0, 4)\}$. Alice dan Bob menyepakati titik $B(0, 1)$ sebagai basis.

\begin{enumerate}
    \item Alice memilih $a = 2$, lalu menghitung kunci publiknya:
    \[ P_A = a \cdot B = 2B = B + B = (1, 3) \quad \rightarrow \text{ misalkan titik Q} \]
    \item Bob memilih $b = 3$, lalu menghitung kunci publiknya:
    \[ P_B = b \cdot B = 3B = B + B + B = 2B + B = (3, 3) \quad \rightarrow \text{ misalkan titik R} \]
    \item Alice mengirimkan $P_A$ kepada Bob, Bob mengirimkan $P_B$ kepada Alice.
    \item Alice menghitung kunci rahasia sbb:
    \[ K_A = a \cdot P_B = 2R = R + R = (0, 4) \]
    \item Bob menghitung kunci rahasia sbb:
    \[ K_B = b \cdot P_A = 3Q = Q + Q + Q = (0, 4) \]
\end{enumerate}

Jadi, sekarang Alice dan Bob sudah berbagi kunci rahasia yang sama, yaitu $(0, 4)$

\vspace{10mm}
\textcolor{red}{* Sumber bahan: Nana Juhana, Implementasi Elliptic Curve Cryptography (ECC) pada proses Pertukaran Kunci Diffie-Hellman dan Skema Enkripsi El Gamal}

\end{document}
